\clearpage

\section{SMC asintótico}

Considérese el sistema de segundo orden
\[
  \dot x_1=x_2,\qquad
  \dot x_2=f(x,t)+u,
\]
con
\[
  |f(x,t)|\leq L,\qquad
  |\dot f(x,t)|\leq L_d.
\]

Para el seguimiento de una referencia \(r(t)\), se define el error
\[
  e=r-x_1.
\]
La superficie de deslizamiento es
\[
  \boxed{\sigma=\dot e+me},
  \qquad m>0.
\]
Como
\[
  \dot e=\dot r-x_2,
\]
se obtiene
\[
  \sigma=\dot r-x_2+m(r-x_1).
\]

Se introduce la superficie auxiliar
\[
  \boxed{S=\dot\sigma+m_s\sigma},
  \qquad m_s>0.
\]
La condición \(S=0\) implica \(\sigma\to0\). Además,
\[
  \sigma=\dot e+me=0
\]
describe una dinámica estable del error, por lo que \(e\to0\).

\subsection{Derivación de la ley de control}

La derivada de la superficie principal es
\[
  \dot\sigma
  =\ddot r-\dot x_2+m\dot e
  =\ddot r-f(x,t)-u+m(\dot r-x_2).
\]

Por tanto,
\[
  \ddot\sigma
  =\dddot r-\dot f(x,t)-\dot u
  +m\bigl[-f(x,t)-u\bigr]
  +m\bigl(\ddot r\bigr).
\]

Derivando \(S\), se tiene
\[
  \dot S=\ddot\sigma+m_s\dot\sigma.
\]
Al agrupar los términos,
\begin{align*}
  \dot S
  &=
  \dddot r+m\ddot r
  +m_s\bigl(\ddot r+m\dot r-mx_2\bigr)\\
  &\quad
  -(m+m_s)\bigl[f(x,t)+u\bigr]
  -\dot f(x,t)-\dot u.
\end{align*}

Se define
\[
  V=\frac{1}{2}S^2.
\]
Para referencias y perturbaciones acotadas, se propone el controlador dinámico
\[
\begin{aligned}
  \dot u={}&
  \dddot r+m\ddot r
  +m_s\bigl(\ddot r+m\dot r-mx_2\bigr)\\
  &-(m+m_s)u
  +K\operatorname{sign}(S).
\end{aligned}
\]
Con esta elección,
\[
  \dot S
  =-K\operatorname{sign}(S)
  -(m+m_s)f(x,t)-\dot f(x,t).
\]
Así,
\[
  \dot V
  \leq
  -\left[K-L_d-(m+m_s)L\right]|S|.
\]
Por consiguiente, una condición suficiente de estabilidad es
\[
  \boxed{K>L_d+(m+m_s)L}.
\]

\subsection{Aplicación al modelo PVTOL}

Después de la linealización exacta por realimentación, el modelo desacoplado
es
\[
  \ddot x=v_x,\qquad
  \ddot z=v_z,\qquad
  \ddot\theta=\tau.
\]

Para cada canal \(q\in\{x,z,\theta\}\), se define
\[
  e_q=r_q-q,
  \qquad
  \sigma_q=\dot e_q+m_qe_q.
\]
En particular,
\begin{align*}
  \sigma_x&=\dot r_x-\dot x+m_x(r_x-x),\\
  \sigma_z&=\dot r_z-\dot z+m_z(r_z-z),\\
  \sigma_\theta&=\dot r_\theta-\dot\theta
  +m_\theta(r_\theta-\theta).
\end{align*}

Las superficies auxiliares son
\begin{align*}
  S_x&=\dot\sigma_x+m_{sx}\sigma_x,\\
  S_z&=\dot\sigma_z+m_{sz}\sigma_z,\\
  S_\theta&=\dot\sigma_\theta+m_{s\theta}\sigma_\theta.
\end{align*}

Las entradas \(v_x\), \(v_z\) y \(\tau\) se consideran estados de control
dinámicos. En la implementación considerada se emplean referencias
constantes para cada canal, por lo que
\[
  \dot r_q=\ddot r_q=\dddot r_q=0,
  \qquad q\in\{x,z,\theta\}.
\]
En consecuencia, sus derivadas se calculan mediante
\begin{align*}
  \dot v_x={}&
  -m_xm_{sx}\dot x
  &-(m_x+m_{sx})v_x
  +K_x\operatorname{sign}(S_x),\\
  \dot v_z={}&
  -m_zm_{sz}\dot z
  &-(m_z+m_{sz})v_z
  +K_z\operatorname{sign}(S_z),\\
  \dot\tau={}&
  -m_\theta m_{s\theta}\dot\theta
  &-(m_\theta+m_{s\theta})\tau
  +K_\theta\operatorname{sign}(S_\theta).
\end{align*}

Las condiciones suficientes para los tres canales son
\[
\begin{aligned}
  K_x&>L_{d,x}+(m_x+m_{sx})L_x,\\
  K_z&>L_{d,z}+(m_z+m_{sz})L_z,\\
  K_\theta&>L_{d,\theta}
  +(m_\theta+m_{s\theta})L_\theta.
\end{aligned}
\]

La referencia angular se obtiene de la linealización:
\[
  r_\theta:=\theta_d
  =\operatorname{atan2}(-v_x,v_z+g).
\]
Finalmente,
\[
  \boxed{U_d=\sqrt{v_x^2+(v_z+g)^2}}.
\]